\begin{otherlanguage}{english}
\chapter*{\centering Abstract}
Inventory management requires recognizing when to adapt replenishment to demand context and when to retain a reference policy. This dissertation proposes the Adaptive Inventory Policy Engine (AIPE), a modular, traceable methodology supporting logistics and inventory managers in reducing avoidable costs without practically important deterioration in service. The procedure comprises five stages: specifying the decision, representing context, building comparable evidence, recommending a policy, and evaluating its consequences. Its Policy Selection Engine (PSE) maps profiles to policy families or estimates recommendations for each series, with an intervention rule that can retain the reference. We evaluate seven policies in two separate applications: 145 business series from Bahia at two origins and a historical cohort of 2,408 identifiers from the fifth Makridakis forecasting competition (M5), containing Walmart sales, at three origins. Cost, service, deterioration, and distance from retrospective references assess the utility of each implementation. In Bahia, profile-based selection adds service at an additional cost; class prediction increases cost and reduces service, while prediction of consequences with fallback retains the reference. In M5, prediction of consequences with fallback reduces mean cost by 1.96\% and adds 0.32 percentage points of service, although individual deterioration occurs and mean service remains below 70\%. Oracles distinguish selection errors, portfolio gaps, and physical infeasibility. The central contribution is the decision-support procedure, supported by integration of the PSE approaches, diagnosis of limitations, and evaluation in complementary contexts. The results exercise its interfaces and delimit the utility of the implementations, without establishing realized financial savings or cycle-by-cycle PSE revision. Historical cohorts, scenario-based costs, short horizons, and prior development restrict generalization and require independent confirmation.

\vspace{2em}
\textbf{Keywords}: Selection methodology; Inventory management; Contextual policy selection; Service level; AIPE.
\end{otherlanguage}
